\textbf{Job-specific technical preparation.}
To make my fit inspectable rather than merely asserted, I developed
\href{https://github.com/pkd-45/radio-signal-intelligence-lab}{\texttt{radio-signal-intelligence-lab}},
a compact, tested pipeline that converts paired ON/OFF radio dynamic spectra into a balanced, provenance-tracked dataset, trains a purpose-built multimodal PyTorch CNN, exports the model with TorchScript, and produces a transparent candidate-ranking table. The model combines waterfall images with interpretable signal metadata and is deliberately small enough to train and test reproducibly on a laptop. An optional \texttt{blimpy} adapter provides a path to Breakthrough Listen filterbank and HDF5 products, while the public release keeps synthetic benchmark results clearly separate from real-data claims. This complements my existing \texttt{ml-radio-recon} and \texttt{radio-pipeline-control-lab} repositories, which demonstrate radio-interferometric deep learning, anomaly detection, tested pipeline execution, restart/resume, integrity checks, telemetry, and Slurm planning.
